% -----------------------------------------------------------------------------------
\section{Philosophical discussion}\label{sec:discussion}

The preceding sections prove theorems about stated definitions. This section interprets them. Its
claims are arguments, not results, and where they go beyond what has been proved they are marked as
such. Its aim is the one that motivated the inquiry: to say exactly where the disagreement among
decision theories is located, so that the apparent paradoxes can be separated from the substantive
question they conceal.

\subsection{What a formalisation decides before any theorem is proved}\label{sec:disc-format}

The first lesson concerns method, and the paper is its own evidence. Part~\ref{part:single}
compared decision theories on problems that supply, for every policy $\pi$, the law $\mu_\pi$ of the
world's response (Definition~\ref{def:structural}). On that class, ex-ante optimality is attained by
the most naive construction available, choosing a maximiser problem by problem
(Theorem~\ref{thm:universal}), and every impossibility for criteria that judge each problem
separately reduces to a single problem on which two acceptable sets are disjoint
(Lemma~\ref{lem:pointwise}). The result is cheap for the reason Section~\ref{sec:programs} makes
explicit: causal, evidential, and functional decision theory disagree precisely about where $\mu_\pi$
comes from, and a model that supplies $\mu_\pi$ has already settled their dispute. The
well-definedness requirement of Definition~\ref{def:welldefined} did the same for problems that refer
to a particular theory. The fairness condition admitted them, since fixed algorithms may occur as
constants (Definition~\ref{def:fair}), but Remark~\ref{rem:skalse} found the corresponding
construction ill defined, because it needed a fixed point that no well-defined finite problem
supplies. The verdict was reversed once problems became programs (Theorem~\ref{thm:fixedcopy}).

The general point is this. To state a decision problem formally, one must say how the agent's
alternatives are to be evaluated, and that choice is usually the one in dispute. A formal result
about decision theories is therefore only as informative as the specification of alternatives it
presupposes, and the specification is often made in the format of the problems rather than in the
statement of the theorem, where it is easy to miss. Appendix~\ref{app:stipulations} and
Section~\ref{sec:requirements} exist to prevent this. Reading the literature with the same question
in mind, namely \emph{which part of the answer was built into the format?}, removes much of the
appearance that one theory has been refuted by a problem rather than by an assumption.

\subsection{The locus of the conflict}\label{sec:disc-locus}

The theories compared here agree about more than their dispute suggests. They share the payoffs,
the probabilities of the world's own randomness, and the rule of maximising expected value.
Section~\ref{sec:model} isolates what remains in a single subscript. The causal value
$Q_\pi(h,a)$ of an alternative $a$ at $h$ averages with $\mu_\pi$, the law induced by the policy
actually adopted; the ex-ante value $V_D(\dev{\pi}{h}{a})$ averages with $\mu_{\dev{\pi}{h}{a}}$, the
law the alternative itself would induce. Every within-problem conflict of Part~\ref{part:single} is a
consequence of that difference. For a single choice, Theorem~\ref{thm:characterisation} states
exactly when it can arise: when, and only when, the law of the states preceding the choice depends on
what the agent prescribes. Proposition~\ref{prop:causaltree} gives the complementary fact: if the
world's response does not depend on the agent's policy, every ex-ante optimum is causally optimal
and no conflict arises.

The same choice reappears in each part of the paper, in a different form.
\begin{itemize}
  \item In Part~\ref{part:single} it is the choice of the law under which an alternative is
        evaluated.
  \item In Part~\ref{part:across} it is the choice between the row and the diagonal of the causal
        matrix, and the cross-problem results depend on whether the belief perspectives associated
        with removed options disappear with them (Section~\ref{sec:menus}).
  \item In Part~\ref{part:programs} oracle substitution recovers the diagonal and port-only
        intervention a row (Proposition~\ref{prop:progmatrix}). At the axiomatic level the choice is
        one between two notions of when computations are the same
        (Definition~\ref{def:equivalence}): grounded equivalence, which identifies components by
        their behaviour once the actual agent is inserted, and parametric equivalence, which
        identifies them only if they agree under every substitution. Calibration together with
        grounded equivalence forces causal evaluation (Lemma~\ref{lem:projection}); parametric
        equivalence admits ex-ante evaluation satisfying every desideratum
        (Corollary~\ref{cor:parametric}).
\end{itemize}

This, then, is the locus of the conflict. The theories do not disagree about how to choose among
specified alternatives. They disagree about what the alternatives to the agent's own choice
\emph{are}. Causal decision theory evaluates an alternative in a world that keeps the actual past,
including whatever in it predicted the agent, and varies only the act: the small miracle of
\textcite[468]{lewis1979arrow}, located at the moment of choice. Functional decision theory evaluates an
alternative in a world in which the agent's decision procedure outputs something else, so that
everything computing that procedure varies with it. Neither world is actual, and nothing in the
actual course of events selects one of them. The discussion following
Proposition~\ref{prop:doublecall} states this in its simplest form: a fixed program computing $0$ and
a linked call to an agent that outputs $0$ agree on the actual execution and differ under
replacement. Theorem~\ref{thm:transport} states it in general: no total semantics can both transport
a computation through alternative outputs and identify every presentation with the same actual
behaviour, because actual behaviour discards the information that transport needs.

The philosophical literature on counterpossibles meets the same obstruction from the other side
(Section~\ref{sec:cp}). Williamson's argument for vacuism identifies antecedents by their
possible-worlds intension, which is empty for every impossibility; grounded extensionality identifies
components by their instantiated behaviour, which is the same for the identity context and the
constant context at the baseline. In both cases an identification appropriate to what actually
happens collapses the distinctions that non-vacuous counterfactuals need.
\textcite{wolpertbenford2013} reached a related diagnosis of Newcomb's problem: the two classical
arguments presuppose incompatible formalisations of how choice and prediction are related. The
present results locate that incompatibility and show that it recurs, in the same form, wherever the
agent's alternatives must be specified.

\subsection{Not arbitrary: the coupling between agent and predictor}\label{sec:disc-coupling}

If the facts of the actual execution do not decide between the two evaluations, one might conclude
that the choice is a convention, and that the dispute between causal and functional decision theory
is verbal. The results do not support that conclusion, for two reasons.

First, the choice has consequences. Agents built on the two evaluations, placed before the same
reliable predictor, obtain different payoffs (Proposition~\ref{prop:newcomb}). A dispute whose
resolution changes what agents receive is not merely verbal.

Second, the candidate ways of making the choice are constrained. Section~\ref{sec:inference} shows
that the natural rules for recognising the agent in a closed program each fail in a characteristic
way: literal matching misses encoded copies; local matching manufactures dependence from accidental
agreement (Theorem~\ref{thm:localmatch}); global matching misses transformed dependence such as $1-T$
(Theorem~\ref{thm:complement}) and can make optimality decide the halting problem
(Theorem~\ref{thm:haltred}); and implication-based reasoning about the actual agent yields spurious
values (Theorem~\ref{thm:fiveten}) or none (Proposition~\ref{prop:guardedarith}). A semantics is not
something one may simply adopt; it must survive these failures.

What, then, could fix the choice? The proposal of this section, which goes beyond anything proved
above, is that it is fixed, where it is fixed at all, by the \emph{mechanism} through which the
predictor's output depends on the agent. The actual execution does not reveal that mechanism, but
the mechanism is a feature of the world, and in some cases it is known.
\begin{itemize}
  \item A predictor that runs the agent's own code on the agent's own input, through what
        Section~\ref{sec:programs} calls a designated call, depends on the agent's decision function
        in the parametric sense. Changing the function's output changes the prediction, and oracle
        substitution, hence ex-ante evaluation, describes the situation literally
        (Theorem~\ref{thm:promiseopt}).
  \item A predictor whose output is fixed by data that do not compute the agent's function, such as
        base rates in a population or a record made after the fact, does not depend on that function
        at all. Port-only intervention describes it correctly, and the two evaluations coincide
        whenever the induced laws do (Proposition~\ref{prop:causaltree}).
  \item The difficult cases lie between, and they are the ones thought experiments typically
        present. A predictor that scans the agent's brain beforehand reads a physical state that will
        later produce the decision. Whether its output depends on the agent's decision function, or
        only on a physical state with which the decision is correlated, depends on whether the agent
        is identified with the computation the brain implements. That identification is not a fact
        about the prediction mechanism alone; it is a commitment about what an agent is.
\end{itemize}

On this view, the error of the debate has been to seek a decision rule that is correct independently
of the coupling between agent and world. No such rule should be expected. What can be asked of a
theory is a correct evaluation \emph{given} the coupling, and causal and ex-ante evaluation each
supply one, for different couplings. The irreducibly philosophical residue is then smaller and better
defined than the debate suggests. It is not \enquote{which decision theory is rational?}, but whether,
in the cases where the coupling runs through an agent's physical constitution, the agent is that
constitution or the function it computes. The extraction problem left open in
Section~\ref{sec:conclusion}, whether a dependency graph can be recovered from an arbitrary closed
program, is the formal shadow of this question, and Corollary~\ref{cor:parametric} shows that once it
is answered nothing further stands in the way of a coherent semantics.

\subsection{When evaluation happens: the Bomb and updatelessness}\label{sec:disc-bomb}

The Bomb is the case most often offered against functional decision theory, and its treatment in
Section~\ref{sec:bomb} shows the same locus at work. As usually stated, the problem is not well
defined (Proposition~\ref{prop:bombill}): a predictor whose forecast determines the observation, and
whose target is the agent's eventual action, cannot be accurate for every policy. Under the policy
$(L,R)$, which takes Left after seeing a bomb and Right otherwise, every forecast is wrong. Under the
policy $(R,L)$, every forecast is right: a forecast of Right produces a bomb and the agent takes
Right, while a forecast of Left produces no bomb and the agent takes Left. The problem therefore
leaves undetermined what such an agent receives. The repaired version
(Proposition~\ref{prop:bomb}) fixes the predictor's target as the agent's prescription on seeing a
bomb, and under that stipulation ex-ante evaluation selects $(L,L)$ while causal evaluation selects
$(R,L)$. The force of the case thus depends on an assumption that its informal statement does not
make.

Even repaired, the disagreement concerns the vantage point of evaluation, not the value of outcomes.
Section~\ref{sec:bomb} records both readings. Taken as the policy problem, evaluated before the
prediction, the Bomb favours Left: a risk of death of probability $\varepsilon=10^{-24}$ against a
near-certain cost of $1$. Taken as the already-conditioned problem, \enquote{a bomb is visibly
present}, ex-ante evaluation itself selects Right, since the options are then worth $-M$ and $-1$.
Functional decision theory evaluates from before the prediction and declines to update on what it
sees, which is what \emph{updatelessness} means \autocite{dai2009,yudkowskysoares2017}. The objection
to it is that the problem posed is the second one, so that a calculation conducted from the first
answers a different question. Counterfactual mugging (Section~\ref{sec:mugging}) separates the same
two vantage points with a prediction about an information set that is unreached on half the paths.

The policy vantage point is not unreasonable. The literature on resolute choice and dynamic
consistency \autocite{mcclennen1990,gauthier1994} and the planning-stage analysis of the absent-minded
driver \autocite{piccionerubinstein1997,aumannetal1997} defend executing an ex-ante plan even where it is
not optimal ex post. And a commitment made before the prediction is worth exactly the ex-ante value of
the policy committed to (Proposition~\ref{prop:commitvalue}), so even a causal decision theorist,
offered the chance to commit at that stage, would commit to Left. The ordinary acceptance of small
risks of death, as in daily travel, makes the policy calculation look sensible. The analogy fails,
however, at exactly the point in dispute: in ordinary risk the bad outcome is not observed at the
moment of choice, whereas in the Bomb the agent chooses after the improbable event has occurred and
is in view. Whether the agent's choice now \emph{is} the output the predictor read earlier, so that
choosing Right now amounts to having been an agent for whom bombs are placed, is the coupling question
of Section~\ref{sec:disc-coupling} in temporal form. If it is, the policy vantage point is correct and
the bomb is the price of a good policy. If it is not, the agent simply stands before a bomb, and the
policy calculation describes a choice that is not being made. The Bomb is therefore neither a
refutation of functional decision theory nor a mere appeal to emotion; it is a test of whether one
accepts the identification on which that theory rests.

\subsection{Fair and unfair problems}\label{sec:disc-fair}

Functional decision theory is claimed to do well on fair problems, those in which the outcome depends
only on the agent's behaviour, and problems built against a particular theory are set aside as
unfair \autocite{yudkowskysoares2017}. The results bear on this restriction in three ways.

First, unfairness is a matter of evaluation semantics, not of the problem's actual behaviour. The
fixed-copy diagonal of Theorem~\ref{thm:fixedcopy} defeats every globally total computable agent when
the embedded copy is held fixed; if the copy is declared linked, the payoff becomes
$\one[A(d)\ne A(d)]=0$ and every agent is optimal. Since the two declarations agree on every actual
execution, the boundary between fair and unfair problems cannot be drawn from behaviour; it is the
linked--fixed distinction of Section~\ref{sec:disc-locus} again. The fairness criterion presupposes a
classification it does not supply, as \textcite[§9, n.~15]{yudkowskysoares2017} acknowledge when they note
that it remains unclear how to transform a problem unfair to one theory into an analogous problem for
another. Defences that remove counterexamples by declaring them unfair, or by declaring them outside a
theory's intended domain, do the work that the well-definedness requirement did by stipulation in
Part~\ref{part:single}.

Second, the obstacle to an optimal theory is not choosing well but being computable by the
environments. An optimal decision function exists on both admissible classes, fixed copies of
arbitrary programs included (Theorem~\ref{thm:existence}). What fails is an optimal agent that halts
on every input, because any such agent can be run by an environment that rewards the opposite
(Theorem~\ref{thm:fixedcopy}); the oracles that compute a globally total optimal agent are exactly
the PA degrees
(Theorem~\ref{thm:pa}). Optimal choice is always possible problem by problem; no decision procedure
that the world can run is optimal everywhere. This vindicates the intuition that a counterexample to
any predefined theory is always easy to construct, and it locates that intuition precisely: the
counterexample needs the theory to be a program the environment can execute.

Third, the intuition should not be overstated. Problems built against a theory exist for every
theory the environments can compute, but not only for those: a non-computable decision function can
also be predictably wrong (Proposition~\ref{prop:unfair}). What protects an optimal function is
optimality, not non-computability. The slogan that unfair problems exist exactly for computable
theories is false; the defensible statement is that global simulability by the environments
guarantees a constructible defeat.

\subsection{Thought experiments as instruments}\label{sec:disc-thought}

Idealisation is not the difficulty with these thought experiments. Physics proceeds through
frictionless planes, and a reliable predictor is a legitimate device for isolating one feature of a
decision, its predictability. The difficulty is that the device is usually specified in exactly the
respect that matters least and left unspecified in the respect that matters most. A problem states
the predictor's accuracy, sometimes to twenty-four decimal places, and says nothing about its
mechanism, its target, or how its errors arise. The results above show that these determine the
verdicts.
\begin{itemize}
  \item The \emph{mechanism} fixes whether the prediction is linked to the agent
        (Sections~\ref{sec:disc-locus} and~\ref{sec:disc-coupling}).
  \item The \emph{target} decides whether the Bomb is even well defined
        (Proposition~\ref{prop:bombill}), and restrictions on admissible targets decide which
        conflicts survive (Section~\ref{sec:targets}).
  \item The \emph{error kernel} is a stipulation (item~\ref{stip:error} of
        Appendix~\ref{app:stipulations}). \enquote{Wrong once in $10^{24}$ cases} is an overall rate.
        Whether errors are independent of the agent's policy, concentrated in particular situations,
        or correlated with the agent's disposition changes the ex-ante calculation, and no statement
        of the problem fixes it.
\end{itemize}

Two further habits compound the underspecification. Extreme magnitudes, such as an error rate of
$10^{-24}$ against a death cost of $10^{9}$, push intuitions into a range in which they have no
evidential basis. And informal arguments import structure that the problem does not contain: repeated
play, where the problem is one-shot; or, in the other direction, the stipulation that the agent knows
it is not inside the predictor's simulation, which the problem may not license.
Section~\ref{sec:requirements} lists the assumptions that any optimality theorem for an agent reading
descriptions must state. The informal debate would benefit from the corresponding discipline: no
verdict on a predictor problem should be accepted until the predictor's mechanism, target, and error
kernel are fixed, because until they are, the problem has no determinate answer to dispute.

\subsection{What the results leave untouched: the utility presupposition}\label{sec:disc-utility}

Every result in this paper assumes that outcomes have real-valued payoffs
(Definition~\ref{def:problem}). Causal, evidential, and functional decision theory share this
assumption, and the critique it invites applies to all of them equally.

In the strict sense required by the representation theorems of \textcite{vonneumannmorgenstern1944}
and \textcite{savage1954}, an agent has a utility function when a single, menu-independent ranking of
outcomes, complete, transitive, and continuous, accounts for its choices in every situation. Agents
with incomparable or lexicographic values fail this. Consider a forced choice between two children,
the predicament of \textcite{styron1979}, and model the chooser as a coalition of two sub-agents, each
protecting one child at any price. A small bonus attached to either option does not tip the choice,
so the options are incomparable rather than tied, which is the small-improvement argument of
\textcite{chang2002}; and a lexicographic preference has no real-valued representation at all
\autocite[159, with n.~1 on p.~164]{debreu1954}. If the coalition responds to deadlock by refusing, it brings about the outcome
both sub-agents rank last. The menu dependence such an agent exhibits violates exactly the contraction
condition \ax{C} of Part~\ref{part:across}, Sen's property~$\alpha$: an option chosen from a larger
menu is not chosen from a smaller one that contains it. Harsanyi's aggregation theorem
\autocite[§III]{harsanyi1955} adds that if the coalition and its members were all expected-utility
maximisers and the coalition respected their unanimous preferences, its utility would be a weighted
sum of theirs, so it would trade one commitment against the other, which is precisely what such an
agent refuses to do.

There is a weaker sense in which every consistent agent maximises something: indexing outcomes by
the situation in which they occur, or redescribing them finely enough, always yields a function that
the agent's choices maximise \autocite[ch.~5]{broome1991}. That sense is always available and therefore says
nothing. It is the counterpart, for utilities, of the lookup tables of Theorem~\ref{thm:universal}.
Enriching outcomes with factors such as guilt or responsibility can restore a strict representation,
sometimes rightly, since the obvious ranking of outcomes may not be the agent's actual one. But each
such enrichment must make predictions across situations, or it becomes the vacuous representation
again.

Two consequences follow for the scope of the paper. Its results are conditional on agents with
well-defined payoffs, and they do not show that real agents have them. And the epistemic objection
to applying decision theory, that an agent's utilities are opaque, shaped by culture and unavailable
to introspection, is distinct from the metaphysical one, that for some agents there is no utility
function to be known. The first limits application; the second limits the subject matter of the
theory.

\subsection{Practical scope}\label{sec:disc-practice}

Where the world's response does not depend on the agent's policy, all the evaluations considered here
coincide (Proposition~\ref{prop:causaltree}). That covers most practical decisions, and in that domain
the shared core, expected-value reasoning under uncertainty, has a long record of use. The distinctive
content of the theories is confined to cases of policy dependence, and those are the cases in which
each fails some criterion: causal evaluation cannot always be ratified (Proposition~\ref{prop:did}),
ex-ante evaluation acts against what the agent observes (Proposition~\ref{prop:bomb}), and every agent
that halts on all inputs can be diagonalised by an environment that runs it
(Theorem~\ref{thm:fixedcopy}). The following is an argument, not a result: a theory whose value lies
only in the cases where its rivals disagree, and which fails some reasonable criterion in those same
cases, has no demonstrated practical advantage over its rivals. Nothing in this paper suggests
otherwise.

The one setting in which the coupling view of Section~\ref{sec:disc-coupling} promises more is that
of agents whose code can be inspected or run by others. There the mechanism of prediction can be
known, the dependence can literally be a designated call, and Theorem~\ref{thm:promiseopt} gives a
computable optimiser for it. Even there, Theorem~\ref{thm:fixedcopy} applies whenever the environment
can run the agent with its output held fixed. For human agents facing hypothetical predictors of
unspecified mechanism, the dispute between causal and functional decision theory has, on the view
defended here, no determinate answer, because the information that would decide it has not been
given.

\subsection{Summary}

The disagreement between causal and functional decision theory is not a disagreement about how to
choose. It is a disagreement about what the alternatives to one's own choice are. Formally, it is the
choice between the row and the diagonal of the causal matrix (Proposition~\ref{prop:progmatrix});
semantically, between grounded and parametric equivalence of computations
(Definition~\ref{def:equivalence}). The actual course of events does not settle that choice
(Theorem~\ref{thm:transport}), and the choice is not arbitrary (Section~\ref{sec:inference}). This
section has proposed that it is settled, where it is settled at all, by the mechanism coupling the
agent to whatever depends on it, with an irreducible residue concerning whether an agent is identified
with the computation it implements. Where a thought experiment leaves that mechanism unspecified, the
dispute has no determinate answer, and the appearance of paradox is what that underspecification
leaves behind.
